%% kapa: exam=2 round=32 year=2021 period=3 subject=compensation_law subject_ko=감정평가 및 보상법규 seq=1 number=1 points=40
%% source: Q-Net 감정평가사 기출문제 (cst00309) artlSeq=5208993 file=raw/1008/5208993/감평 2차 시험지(2021).pdf
%% conversion: kordoc 4.24.51 -> tools/gen2.py (auto); review=proofread-sonnet; flags=none
\begin{longquestion}{【문제 1】}
국토교통부장관은 2013. 11. 18. 사업시행자를 ‘A공사’로, 사업시행지를 ‘X시 일대 8,958,000 m\textsuperscript{2}’로, 사업시행기간을 ‘2013. 11.부터 2017. 12.까지’로 하는 ‘◇◇공구사업’에 대해서 「공익사업을 위한 토지 등의 취득 및 보상에 관한 법률」에 따른 사업인정을 고시하였고, 사업시행기간은 이후 ‘2020. 12.까지’로 연장되었다. 甲은 ㉮토지 78,373 m\textsuperscript{2}와 ㉯토지 2,334 m\textsuperscript{2}를 소유하고 있는데, ㉮토지의 전부와 ㉯토지의 일부가 사업시행지에 포함되어 있다. 종래 甲은 ㉮토지에서 하우스 딸기농사를 지어 왔고, ㉯토지에서는 농작물직거래판매장을 운영하여 왔다. 甲과 A공사는 사업시행지 내의 토지에 대해 「공익사업을 위한 토지 등의 취득 및 보상에 관한 법률」에 따른 협의매수를 하기 위한 협의를 시작하였다. 다음 물음에 답하시오. (아래의 물음은 각 별개의 상황임) (40점) \par
\medskip\noindent 물음 1) 협의 과정에서 일부 지장물에 관하여 협의가 이루어지지 않아 甲이 A공사에게 재결신청을 청구했으나 A공사가 재결신청을 하지 않는 경우, 甲의 불복 방법에 관하여 검토하시오. (15점) \par
\medskip\noindent 물음 2) ㉮토지에 대하여 협의가 성립되지 않았고, A공사의 수용재결신청에 의하여 ㉮토지가 수용되었다. 甲은 ㉮토지가 수용되었음을 이유로 A공사를 상대로 「공익사업을 위한 토지 등의 취득 및 보상에 관한 법률」에 따른 재결 절차를 거치지 않은 채 곧바로 농업손실보상을 청구할 수 있는지를 검토하시오. (10점) \par
\medskip\noindent 물음 3) 협의가 성립되지 않아 사업시행지 내의 ㉯토지가 수용되었다. 그 후 甲은 ㉯토지의 잔여지에 대해서 2020. 11. 12. 잔여지수용청구를 하였다. 잔여지수용청구권의 법적 성질과 甲의 잔여지수용청구가 인정될 수 있는지를 검토하시오. (15점) \par
\medskip
\begin{framedtext}
\textbf{참조 조문} \par
\textbf{<공익사업을 위한 토지 등의 취득 및 보상에 관한 법률>} \par
\textbf{제28조(재결의 신청)} ① 제26조에 따른 협의가 성립되지 아니하거나 협의를 할 수 없을 때(제26조제2항 단서에 따른 협의 요구가 없을 때를 포함한다)에는 사업시행자는 사업인정고시가 된 날부터 1년 이내에 대통령령으로 정하는 바에 따라 관할 토지수용위원회에 재결을 신청할 수 있다. \par
② <생략> \par
\textbf{제30조(재결 신청의 청구)} ① 사업인정고시가 된 후 협의가 성립되지 아니하였을 때에는 토지소유자와 관계인은 대통령령으로 정하는 바에 따라 서면으로 사업시행자에게 재결을 신청할 것을 청구할 수 있다. \par
② <이하 생략> \par
\textbf{제77조(영업의 손실 등에 대한 보상)} ① 영업을 폐업하거나 휴업함에 따른 영업손실에 대하여는 영업이익과 시설의 이전비용 등을 고려하여 보상하여야 한다. ② 농업의 손실에 대하여는 농지의 단위면적당 소득 등을 고려하여 실제 경작자에게 보상하여야 한다. 다만, 농지소유자가 해당 지역에 거주하는 농민인 경우에는 농지소유자와 실제 경작자가 협의하는 바에 따라 보상할 수 있다. \par
③ <이하 생략> \par
\end{framedtext}
\end{longquestion}
